% =====================================================================
%  PROPUESTA DE PILOTAJE (para presentar a la profesora Teresa Pontón)
%  Resume qué se implementa, cómo se adaptaron sus ejercicios y qué se
%  le pide. No incluye el análisis a priori (está en secuencia_tareas).
%  Compilar con: pdflatex propuesta_pilotaje.tex (dos veces)
% =====================================================================
\documentclass[11pt,letterpaper]{article}

\usepackage[utf8]{inputenc}
\usepackage[T1]{fontenc}
\usepackage[spanish,es-tabla,es-nodecimaldot]{babel}
\usepackage{lmodern}
\usepackage{microtype}
\usepackage[margin=2.2cm]{geometry}
\usepackage{amsmath,amssymb}
\usepackage{booktabs,tabularx,array}
\usepackage{enumitem}
\usepackage[table]{xcolor}
\usepackage[most]{tcolorbox}
\usepackage{fancyhdr}
\usepackage{lastpage}
\usepackage{setspace}

\definecolor{unalverde}{RGB}{0,98,65}
\definecolor{grisclaro}{RGB}{242,242,242}
\DeclareMathOperator{\seno}{sen}
\newcolumntype{L}{>{\raggedright\arraybackslash}X}
\setlist{itemsep=2pt,topsep=3pt}
\setstretch{1.08}

\newtcolorbox{caja}[1]{colback=grisclaro,colframe=unalverde,fonttitle=\bfseries,title={#1},breakable,arc=2pt}

\pagestyle{fancy}
\fancyhf{}
\fancyhead[L]{\small Propuesta de pilotaje --- Secuencia de tareas con Gemini}
\fancyhead[R]{\small UNAL -- Sede Palmira}
\fancyfoot[C]{\small Página \thepage\ de \pageref{LastPage}}
\renewcommand{\headrulewidth}{0.4pt}

\begin{document}

{\LARGE\bfseries Propuesta de pilotaje\par}
\smallskip
{\large Secuencia de tareas de cálculo diferencial con el acompañamiento de Gemini\par}
\medskip

\begin{tabularx}{\linewidth}{@{}>{\bfseries}p{2.4cm}L@{}}
Para & Prof. Teresa Pontón Ladino, Cálculo Diferencial \\
De & Dairo Janover Peña Cárdenas, tutor del programa GEA y estudiante de la Maestría en Enseñanza de las Ciencias Exactas y Naturales \\
Directores & Dra. Marisol Santacruz Rodríguez y Dr. Luis Octavio González Salcedo \\
Fecha & \rule{4cm}{0.4pt} \\
\end{tabularx}

% =====================================================================
\section*{1. De qué se trata}

Mi trabajo de grado estudia cómo resuelven problemas de cálculo los estudiantes remitidos a tutoría personalizada cuando trabajan primero solos y luego con Gemini, la IA generativa de la cuenta institucional. No busco medir notas. Me interesa ver qué cambia en sus estrategias, sus errores y sus justificaciones cuando la IA entra a acompañar el trabajo, y qué condiciones de la tarea evitan que la IA les resuelva el ejercicio.

Para eso armé una secuencia de cuatro tareas con ejercicios de los talleres de clase. Tres vienen de su \emph{Taller Módulo 1} (ejercicios 24 y 30) y uno del taller de la profesora Mónica Mosquera (ejercicio 8). Antes de aplicarla con los estudiantes del estudio necesito probarla con uno o dos estudiantes que no participen en él. Por eso le presento la propuesta.

% =====================================================================
\section*{2. Cómo se adaptaron sus ejercicios}

El problema matemático se conserva. Lo que cambia es lo que se le pide al estudiante: en lugar de ``halle el límite'', se le pide juzgar razonamientos, comparar representaciones o decidir sin que se le indique el camino.

\begin{center}\footnotesize
\renewcommand{\arraystretch}{1.25}
\begin{tabularx}{\linewidth}{@{}p{3.3cm}LL@{}}
\toprule
\textbf{Ejercicio del taller} & \textbf{Qué se pide en la tarea} & \textbf{Qué se espera observar} \\
\midrule
\textbf{T1.} Ej.~24.6\newline $\displaystyle\lim_{x\to3}\frac{x^4-81}{x^2-x-6}$ & Juzgar cuatro razonamientos de otros estudiantes: ``$\frac00$, no existe''; una factorización con el error $(x^2-9)^2$; la correcta; y ``la calculadora da 21,54''. Precisar el dominio. & Con qué criterio valida el estudiante: un argumento, un cálculo o la autoridad de quien lo dice. Si detecta el error de álgebra previa. \\
\midrule
\textbf{T2.} Ej.~30.7\newline $\displaystyle\lim_{x\to\infty}\frac{4-x^3}{x^3-1}$ & Discutir la afirmación ``$x^3$ se cancela, el límite es 1''; comparar tabla, álgebra y gráfica; responder la pregunta ``¿asíntotas?'' del propio taller. & Si resuelve el conflicto entre su primera intuición y la tabla. Si encuentra la asíntota vertical. \\
\midrule
\textbf{T3.} Mosquera, ej.~8d\newline $\displaystyle\lim_{x\to0}\frac{\seno 3x+\seno 5x}{2x}$ & Decidir el límite con dos representaciones; comparar radianes y grados; juzgar el atajo ``$\seno 3x+\seno 5x=\seno 8x$'', que da el número correcto con un razonamiento inválido. & Si valida por el resultado o por los pasos. Si conecta el resultado en grados con la conversión de unidades. \\
\midrule
\textbf{T4.} Ej.~24.16 y 31\newline $\displaystyle f(x)=\frac{x+2}{\sqrt{x^2+5}-3}$ & Sin indicar el punto: encontrar dónde no está definida $f$ ($x=\pm2$), decidir en cuál se ``tapa el hueco'' y qué pasa en el otro, e inventar una función propia. & Qué estrategia elige sin indicación. Si distingue una discontinuidad removible de una asíntota vertical. \\
\bottomrule
\end{tabularx}
\end{center}

El orden va de juzgar razonamientos ajenos a producir los propios: cada tarea retira un apoyo que la anterior daba. También sigue el orden de su taller: primero límites en un punto, luego al infinito, y al final una tarea que integra dominio, continuidad y límite, como pide su ejercicio 31.

% =====================================================================
\section*{3. Cómo es una sesión}

Cada tarea se trabaja en una sesión individual de 55 a 60 minutos, en tres momentos:

\begin{enumerate}
  \item \textbf{Sin IA (unos 20 min).} El estudiante resuelve solo, con esfero azul, diciendo en voz alta lo que piensa. El tutor no explica; solo hace preguntas como ``¿qué ha intentado?'' o ``¿cómo sabe que eso es cierto?''.
  \item \textbf{Con Gemini (unos 20 min).} El estudiante consulta Gemini en \emph{modo de aprendizaje guiado}, que responde con preguntas y pistas en lugar de dar la solución. Lo que cambie lo escribe con esfero rojo, sin borrar lo azul. Así la hoja muestra qué cambió después de usar la IA.
  \item \textbf{Validación (unos 10 min).} El estudiante decide su respuesta final, la justifica y dice qué parte hizo él y qué parte hizo Gemini.
\end{enumerate}

Gemini se usa en el computador de la tutoría, con mi cuenta institucional; el estudiante no usa su cuenta ni escribe datos personales.

% =====================================================================
\section*{4. En qué consiste el pilotaje}

\begin{caja}{Lo que se propone}
\begin{itemize}
  \item \textbf{Participantes:} uno o dos estudiantes de su curso, voluntarios. No deben hacer parte después del estudio; si alguno es remitido más adelante a tutoría personalizada, sus datos no se incluirían.
  \item \textbf{Duración:} las cuatro tareas, en cuatro sesiones de una hora o en dos sesiones de dos horas, en el horario que acuerden.
  \item \textbf{Qué se registra:} las hojas de trabajo, la conversación con Gemini y, si el estudiante lo autoriza, el audio. El estudiante firma el mismo consentimiento informado del estudio.
  \item \textbf{Qué no se hace:} el pilotaje no afecta las notas y sus resultados no entran al análisis de la tesis.
  \item \textbf{Qué se prueba:} si las consignas se entienden, si alcanza el tiempo, cómo responde Gemini en modo guiado, si la instrucción inicial funciona y si las hojas sirven para registrar los cambios.
\end{itemize}
\end{caja}

% =====================================================================
\section*{5. Lo que le pido}

\begin{enumerate}
  \item \textbf{Su opinión sobre las tareas:} si las consignas son adecuadas para sus estudiantes y si cambiaría algo.
  \item \textbf{Saber si estos ejercicios ya se resolvieron en clase.} Si se resolvieron en el tablero, los estudiantes podrían recordar la respuesta, y conviene escoger otros del mismo taller.
  \item \textbf{Su ayuda para invitar} a uno o dos estudiantes para el pilotaje.
  \item \textbf{Su autorización para usar y citar sus ejercicios} en el trabajo de grado, con referencia a su taller.
\end{enumerate}

\vspace{0.5cm}
\noindent Muchas gracias por su tiempo. Le llevo impresas las hojas de las cuatro tareas, tal como las recibirá el estudiante.

\vspace{0.8cm}
\noindent\rule{6.5cm}{0.4pt}\\
Dairo Janover Peña Cárdenas\\
dapenaca@unal.edu.co

\end{document}
